\section{Period coordinates and the marked components}\label{sec:streams}

The first selection has made every unmarked complementary component a
point.  It remains to show that, at each marked end, the limiting round
disk fills the entire complementary component.  Suppose that part of the
component protrudes beyond that disk.  A sufficiently small real barrier
excludes the protrusion by the line-payment estimate.  If no such barrier
exists, an endpoint law produces a coordinate with additive period one.
Before smoothing, its boundary variation is a probability measure; its
energy and its oscillations at all other ends can be made arbitrarily
small.  These properties give incompatible lower and upper bounds for a
boundary integral, the upper bound coming from Stokes' theorem.

Throughout this section the prior coordinate list is an arbitrary
countable list satisfying (P).  All references to free points use its
current flag set, enlarged when a new coordinate is appended.

\subsection{Reduction to a compact family of simple arcs}

We first retain only the terminal portions of the endpoint paths and
erase their loops.  Restricting the resulting law to a compact family
will give the uniform winding bounds needed for the period coordinate.

\begin{lemma}[Compact simple-arc reduction]\label{lem:compact-arcs}
Let $\mu$ be an endpoint law supplied by
Theorem~\ref{thm:positive-law}, with target $b$ and prescribed collars
$C_j(b)=\partial U_j(b)$.  For every sufficiently deep prescribed collar
$C=\partial U$, there is a probability law $\eta$ supported on a compact
subset of $C([0,1],X)$ consisting of simple arcs $\alpha$ such that
\[
 \alpha(0)\in C,\qquad \alpha(1)=b,\qquad
 \alpha((0,1))\subset U\setminus\{b\}.
\]
Its mean domain occupation has a density $g_\eta$ satisfying
\begin{equation}\label{eq:arc-occupation}
 g_\eta\le 2\ind_{U\cap G}g_\mu
 \quad\text{almost everywhere},
 \qquad
 m_\eta(s)\le 2m_\mu(s)\quad(s\in B\cap\overline U),
\end{equation}
and $m_\eta=0$ outside $\overline U$.  Here $m$ denotes hit
probabilities, rather than visit counts.  In particular, for every
$\eps>0$, every sufficiently deep prescribed collar admits such a law
$\eta$ with
\begin{equation}\label{eq:arc-smallness}
 \|g_\eta\|_2+
 \|m_\eta|_{B\setminus\{b\}}\|_{\ell^2}
 <\eps.
\end{equation}

For almost every tuple defining any fixed finite prior level union $H$,
the expected number of distinct intersections of an arc with $H$ is
finite, and the arcs almost surely have no free dust contacts in $H$.
If the endpoint problem has gates, the starting point additionally
satisfies $f_0(\alpha(0))\in\overline W_2$.
\end{lemma}

\begin{proof}
Cut each path at its last hit of $C$ and retain its terminal part,
reparametrized on $[0,1]$.
The anchor is outside $\overline U$, so such a hit exists.  Its last
hitting time is strictly less than $1$; after that time the path is
strictly inside $U$.  The original path reaches $b$ for the first time
at its endpoint.  Thus its terminal part has distinct endpoints and
contains a simple arc with those endpoints.  The gate condition passes
to the last hitting point.

These operations can be performed measurably up to completion of the
law.  Last hitting times of a fixed closed set are Borel on the present
path family.  In the product of two path spaces, the conditions that
the second path is injective, has prescribed endpoints, and has image
contained in the first path's image are Borel.  For injectivity, use
\[
 \min_{|s-t|\ge 1/n}d_X(\alpha(s),\alpha(t))>0
 \qquad(n\in\N);
\]
image containment is closed for uniform convergence.  The nonempty
sections therefore admit a universally measurable selection by the
Jankov--von Neumann theorem \cite[Theorems~18.1 and 29.9]{Kechris1995}.  The
pushforward of the completed probability measure is an ordinary Borel
probability measure on path space.

The occupation measure of a selected simple arc is its domain image
length measure.  By the image-length observation in
Section~\ref{sec:barriers}, this is dominated by the occupation of the
original terminal path.  There is no arclength atom at its initial
point on $C$.  Consequently the unconditioned selected-arc law has
occupation density at most $\ind_{U\cap G}g_\mu$, and its hit probabilities are
at most the original hit probabilities.  The endpoint and interior
conditions describe a Borel subset of path space.  Inner regularity
provides a compact subset of that set of probability at least $1/2$.
Conditioning on this subset proves Equation~\eqref{eq:arc-occupation}.

The indicators of $U_j(b)\cap G$ decrease pointwise to zero.  The
square-summable function $m_\mu$ has countable support, and the sets
$(U_j(b)\cap B)\setminus\{b\}$ decrease to the empty set.  Absolute
continuity of the two squared integrals therefore proves
Equation~\eqref{eq:arc-smallness}.  Finally, all selected images are
subsets of the original images.  The level-intersection conclusions of
Lemma~\ref{lem:law-levels} survive restriction, selection, and
conditioning, with their expected counts increased by at most the
conditioning factor.
\end{proof}

\subsection{Averaging the index of a ray}

Each simple arc will determine an integer-valued function on the cyclic
cover of the punctured dust disk.  We average these functions to obtain
a continuous period coordinate, then smooth it while preserving all
nonsingular bands.

\begin{lemma}[A stream coordinate]\label{lem:stream}
Let $U$ be a dust disk, let $b\in B\cap U$, and put $C=\partial U$.
Suppose $\eta$ is supported on a compact family of simple arcs from $C$
to $b$, meeting $C$ only at their initial points and $b$ only at their
terminal points.  Assume that its mean domain occupation has an
$L^2$ density $g_\eta$, that $m_\eta\in\ell^2(B)$, and that it has
the generic level-intersection conclusions of
Lemma~\ref{lem:compact-arcs} for the prior list.

There is a continuous, locally Sobolev function $v$ on the cyclic cover
over $G\cap\overline U$, with additive period one about $b$, such that
\begin{equation}\label{eq:stream-energy}
 \|\dd v\|_{2,G\cap U}\le\sqrt2\,\|g_\eta\|_2.
\end{equation}
At every $s\in B\cap U\setminus\{b\}$ its band is bounded and has
width at most $m_\eta(s)$.  After enlarging the flags by
$\{s\in B:m_\eta(s)>0\}$, its free traces satisfy (P) against the
prior list.
If $\beta$ is the probability distribution of the arc starts on the
positively oriented curve $C$, then $\beta$ is atomless and
\begin{equation}\label{eq:stream-boundary-measure}
 \dd(v|_C)=\beta.
\end{equation}
This equality is understood on a lifted boundary interval: the branch
increases by one over a full turn.

For every $\eps>0$ one may replace $v$ by a smooth coordinate $u$ of
the same period such that
\begin{equation}\label{eq:stream-smoothing}
 \|\dd(u-v)\|_{2,G\cap U}<\eps,
 \qquad \sup_C|u-v|<\eps,
\end{equation}
with the same nonsingular bands and hence with (P).  The discrepancy
$u-v$ is single-valued.  Thus $u$ is an admissible coordinate with
mask $\overline U$ and budget
\[
 a_u(s)=m_\eta(s)\quad(s\in B\cap U\setminus\{b\}),
 \qquad a_u(b)=1,
\]
and zero budget off its mask.
\end{lemma}

\begin{proof}
We first construct $v$ and prove its energy estimate and boundary-measure
identity.  We then control its dust bands and free traces, and finally
produce the smooth admissible coordinate $u$.

\smallskip\noindent\emph{Step 1: the lifted rays.}
Choose an orientation-preserving homeomorphism from
$(\overline U,b)$ to the closed unit disk with its center.  Write the
cyclic cover of $\overline U\setminus\{b\}$ as
\[
 \mathcal H=\R\times[0,\infty),
\]
where the first coordinate is angle divided by $2\pi$, the second is
minus logarithmic radius, and the deck transformation is
$(x,y)\mapsto(x+1,y)$.  These are topological coordinates; all analytic
estimates below will instead use physical conformal charts in $G$.

For an arc $\alpha$, choose the lift starting at $(a,0)$ with
$a\in[0,1)$.  Its image $\Gamma_\alpha$ is a proper simple ray with
$y\to\infty$ at its terminal end.  Its integer translates are disjoint:
an intersection of two translates would give a repeated projected
point of the simple arc.  In the one-point compactification of
$\mathcal H$, the translated ray $\Gamma_\alpha+(k,0)$ is a crosscut
from $(a+k,0)$ to the boundary point at infinity.  Let $R_k(\alpha)$
be its right-hand complementary component, the one incident to boundary
points $(x,0)$ with $x>a+k$.
Jordan separation gives
\[
 R_{k+1}(\alpha)\subset R_k(\alpha).
\]
Indeed the next ray lies on the right of the preceding ray because it
is connected, is disjoint from that ray, and starts on its right.
The right side of the next ray connects to boundary points right of
both starts and is therefore contained in the preceding right side.

\smallskip\noindent\emph{Step 2: the integer index and its local bounds.}
Compactness of the arc family has two consequences.  The arcs approach
$b$ uniformly as the parameter tends to $1$; and, for every
$\delta>0$, their portions with parameter in $[0,1-\delta]$ stay a
uniform positive distance from $b$.  The first assertion follows from
equicontinuity and the common terminal endpoint.  The second follows
from compactness and the fact that every member is injective and meets
$b$ only at parameter $1$.

Fix $0\le Y<\infty$.  Uniform approach to $b$ gives $\delta>0$ such that
every arc point with $0\le y\le Y$ occurs before the common time
$1-\delta$.  On this initial interval the arcs stay uniformly away
from $b$, so their projected circular positions are equicontinuous.
Subdivide time into finitely many intervals on each of which all circular positions
move by less than a semicircle.  Starting with $a\in[0,1)$, successive
lifts then give a uniform bound for all lifted angles on this initial
interval.  Thus all unshifted rays in the strip $0\le y\le Y$ lie
in a common bounded horizontal interval.

For a compact subset of $\mathcal H$, join its points vertically to
the boundary.  Sufficiently large positive or negative translates
miss all those vertical segments.  Their side at the boundary is
fixed by the order of the starts, and it remains fixed along the
segments.  Hence, for every $z$ off the translated rays, the set
$\{k\in\Z:z\in R_k(\alpha)\}$ is nonempty and bounded above.  We may
therefore define
\begin{equation}\label{eq:ray-index}
 J_\alpha(z)=\max\{k\in\Z:z\in R_k(\alpha)\}.
\end{equation}
The same argument bounds $|J_\alpha|$ uniformly over the arc family
on each compact subset of $\mathcal H$, wherever it is defined.
The definition also gives
\begin{equation}\label{eq:ray-index-period}
 J_\alpha(z+(1,0))=J_\alpha(z)+1,
 \qquad J_\alpha(x,0)=\lfloor x-a\rfloor
 \quad(x-a\notin\Z).
\end{equation}

\smallskip\noindent\emph{Step 3: measurability and an intersection inequality.}
The lifted arc on every initial compact time interval is a Borel
function of the arc and its chosen starting angle, by path lifting.
Off a ray, membership in its right component is Borel as well.  One
may test it by countable chains of rational open balls disjoint from
the ray and joining the point to the appropriate boundary interval.
Avoidance of a compact test path is Borel by the initial lifts and a
countable exhaustion of $[0,1)$.  Since a compact path disjoint from a
closed proper ray has positive distance from it, these tests detect
all complement connections.  This proves measurability of the
indices.  On the coincidence set, where $z$ lies on a translated ray,
assign an arbitrary value whenever a random variable is required.  We
will take expectations only at points whose coincidence probability is
zero.

Suppose a curve $\delta$ in the cover projects injectively to
$\overline U\setminus\{b\}$ and its endpoints miss the sample rays.
Every ray separating the endpoint indices must meet $\delta$.
Different such meetings project to different points because the
projection of $\delta$ is injective.  Consequently
\begin{equation}\label{eq:index-intersections}
 |J_\alpha(\delta(1))-J_\alpha(\delta(0))|
 \le \#\bigl(\alpha([0,1])\cap\operatorname{proj}\delta([0,1])\bigr).
\end{equation}
The right side counts distinct projected points; the inequality also
holds when that count is infinite.

\smallskip\noindent\emph{Step 4: averaging and continuity.}
The index is now defined, measurable, and locally uniformly bounded.
We use these properties to average it.
Lemma~\ref{lem:point-avoidance} applies to $\eta$, since its arcs
end at $b\in B$ and have an $L^2$ mean occupation density.
Thus every fixed point of $G\cap\overline U$, including a point
of $C$, is missed almost surely.  For each lift $z$ of such a point,
define
\[
 v(z)=\E_\eta J_\alpha(z).
\]
Local uniform domination makes this expectation finite.  For a
convergent sequence in the lifted domain, discard the null coincidence
events at its terms and its limit.  For every remaining sample the
index is locally constant at the limit.  Dominated convergence proves
continuity of $v$, including at $C$.  The first identity in
Equation~\eqref{eq:ray-index-period} gives additive period one.

\smallskip\noindent\emph{Step 5: the differential estimate.}
Work in a physical conformal rectangle compactly contained in $G\cap U$,
with a common lift.  Here $z=x+iy$ denotes a physical conformal
coordinate, and the spherical metric is $\lambda(z)|\dd z|$; these
$x,y$ are unrelated to the topological angle and depth above.
For a horizontal interval $I$ and height $y$ let
$N_\alpha(I,y)$ count the distinct intersections of the sample arc
with $I\times\{y\}$.  Curve coarea, followed by the occupation bound,
gives, for Borel sets $A$ of heights,
\begin{equation}\label{eq:stream-slice-coarea}
 \int_A\E_\eta N_\alpha(I,y)\,\dd y
 \le \int_{I\times A}\lambda(z)g_\eta(z)\,\dd x\,\dd y.
\end{equation}
To check the conformal factor, Euclidean image length is spherical
image length divided by $\lambda$, whereas spherical area is
$\lambda^2\dd x\dd y$.  Projection coarea costs at most Euclidean
image length.  The same estimate holds after restricting counts to
Borel portions of an interval.

Apply Equation~\eqref{eq:stream-slice-coarea} first to closed intervals
with rational endpoints.  Outside one null set of heights, the mean
counting measure on every compact interior interval is bounded by
$\lambda g_\eta\dd x$.  This follows by the countable interval tests
and then measure approximation.  Averaging
Equation~\eqref{eq:index-intersections} shows on those lines that
\[
 |v(x_2,y)-v(x_1,y)|
 \le\int_{x_1}^{x_2}\lambda(x,y)g_\eta(x,y)\,\dd x.
\]
Continuity permits arbitrary endpoints.  Thus $v$ is absolutely
continuous on almost every horizontal line and
$|\partial_xv|\le\lambda g_\eta$ there.  The vertical argument gives
$|\partial_yv|\le\lambda g_\eta$.  The usual characterization by
absolute continuity on lines yields local Sobolev regularity and
\[
 |\dd v|_{\mathrm{sph}}\le\sqrt2\,g_\eta
 \quad\text{almost everywhere}.
\]
An exhaustion by charts proves Equation~\eqref{eq:stream-energy}.

\smallskip\noindent\emph{Step 6: the boundary measure.}
The distribution $\beta$ of the starting point has no atoms by point
avoidance on $C$.  Between two ordered lifted boundary angles, the
expected difference of the floors in
Equation~\eqref{eq:ray-index-period} counts the mean number of start
translates in that interval.  This proves
Equation~\eqref{eq:stream-boundary-measure}.  In particular, for every
continuous real function $\varphi$ on $C$,
\begin{equation}\label{eq:stream-boundary-integral}
 \int_C\varphi\,\dd v
   =\int\varphi(\alpha(0))\,\dd\eta(\alpha).
\end{equation}
The left side is a Stieltjes integral.  Only the orientation of the
topological angular coordinate is used; no smoothness is required.

We have obtained a continuous period coordinate with the required
energy bound and a positive boundary measure of total mass one.  To
make it admissible for transfer, we must still control its dust bands,
verify compatibility, and smooth it.

\smallskip\noindent\emph{Step 7: bands at the dust.}
Fix $s\in B\cap U\setminus\{b\}$ and a lift of a neighborhood of
$s$.  Consider two sequences of domain points tending to that lift.
At all terms of both sequences the sample index is defined almost
surely. For each sample, Step~2 shows that only finitely many translated
rays meet a compact neighborhood of the chosen lift. Shrink to a connected
neighborhood missing those rays that do not pass through the lift; their
side memberships are then fixed. If the projected sample misses $s$,
the eventual index difference is therefore zero. Otherwise simplicity
leaves just one ray through the lift. Only membership in its right-hand
component can vary, so the index has at most two adjacent values. Hence
the upper limit of the absolute index difference is at most one.
The local uniform bound and the reverse Fatou inequality now give
\[
 \limsup |v(z_n)-v(w_n)|\le m_\eta(s).
\]
Local boundedness makes the upper and lower bands finite, and choosing
sequences approaching their extremes proves the asserted width bound.
When $m_\eta(s)=0$, the unique trace is $\E_\eta J_\alpha(s)$.

\smallskip\noindent\emph{Step 8: compatibility.}
Put $\Sigma'=\Sigma\cup\{s\in B:m_\eta(s)>0\}$, a countable
enlargement of the flags.  Fix a generic finite prior level union $H$.
By hypothesis the rule
\[
 \nu(A)=\E_\eta\#\bigl(A\cap H\cap\alpha([0,1])\bigr),
 \qquad A\subset H\text{ Borel},
\]
defines a finite measure.  It gives zero mass to
$H\cap(B\setminus\Sigma')$, since almost every sample avoids this
whole set.  The counting integrals are measurable by compact incidence
tests and finite distinct-point tests in path space; finite expected
total count then defines the measure on all Borel subsets.

Let $K$ be a nondegenerate component of $H$ and work in a small disk
$V$ with closure in $U\setminus\{b\}$ and with a common lift.
For free points $x,y\in K\cap V$ joined by a simple arc
$\delta\subset K\cap V$, their traces are expectations of the
indices.  Equation~\eqref{eq:index-intersections} gives
\begin{equation}\label{eq:stream-trace-control}
 |v(x)-v(y)|\le\nu(\delta).
\end{equation}
Set $F=K\cap V\cap(B\setminus\Sigma')$.  We have $\nu(F)=0$,
and $K$ is locally connected by Lemma~\ref{lem:level-continua}.
Equation~\eqref{eq:stream-trace-control} verifies the arc hypothesis
of Lemma~\ref{lem:trace-null} for the continuum $K$, the relatively
open set $K\cap V$, the measure $\nu|_K$, and the free set $F$.
Hence its real trace image has Lebesgue outer measure zero.
Countably many components $K$ and lifted disks $V$ suffice.
The singular point $b$ is flagged.  Fubini in the level parameters
therefore proves (P) for the new coordinate.  All incidence statements
may use completed measures, as in the earlier level arguments.

\smallskip\noindent\emph{Step 9: smoothing.}
It remains to smooth $v$ with an error that vanishes at every dust
point, so that all nonsingular bands remain unchanged.
Near the smooth curve $C$, the continuous Sobolev branches extend
across the boundary by reflection in smooth collar charts.  Cover
$G\cap\overline U$ by a locally finite family of relatively compact
simply lifted patches and take a smooth partition of unity.  Smooth a
branch on each patch; on other sheets use its integer translates.
The local smoothing discrepancy is single-valued, so partition
interpolation respects the additive period.

Choose a continuous positive tolerance $e$ on the source patches with
\[
 0<e(x)\le\tfrac12\min\{\eps,\dist_X(x,B)\},
 \qquad x\in G\cap\overline U.
\]
It extends positively across the compact collar of $C$.  On the $j$th
patch choose the branch approximation error in values smaller than
the positive infimum of $e$ on the corresponding compact partition
support.  Simultaneously make its Sobolev error, including the
derivatives of the partition function, less than $\eps2^{-j}$.
Local continuity and Sobolev approximation permit both requirements.
Partition interpolation then gives $|u-v|(x)\le e(x)$, while the
differential errors are summable.  The resulting $u$ satisfies
Equation~\eqref{eq:stream-smoothing} and is smooth up to $C$.
For every domain sequence tending in $X$ to a nonsingular dust point,
$e(x)$ tends to zero; hence the upper and lower band limits are
exactly unchanged in a common local lift.  Its free traces are
therefore unchanged.  Since compatibility for a newly appended
coordinate concerns only those traces against earlier level unions,
(P) is preserved.
\end{proof}

\subsection{A boundary estimate for a transferred stream}

The boundary-measure identity will provide a lower bound near a
protrusion.  The following Stokes estimate supplies the upper bound in
a finite completion, using the energy and the oscillations at the
other holes.

\begin{lemma}[Stokes estimate]\label{lem:stream-stokes}
Let $u$ be a smooth period coordinate with mask $\overline U$ and
singularity $b$, and put $C=\partial U$.  Transfer it to a finite
completion retaining $C$ and the marked hole indexed by $b$.  Write
$F$ for the normalized circle-domain map, $u^*$ for the transferred
coordinate in the target, and $D_b$ for the distinguished round hole.
Assume that the masked region lies in $|w|<L$.  Target lengths and areas
are Euclidean; the norm below is over the domain inside $F(C)$ and
outside its circular holes.  For a target hole $T$ inside
$F(C)$, let $s_T$ be the index of its source terminal disk and $r_T$
its radius.  When $s_T\ne b$, let $A_u(T)$ be its transferred boundary
oscillation.  Then
\begin{equation}\label{eq:stream-stokes}
 \left|\int_{F(C)}\dist(w,D_b)\,\dd u^*\right|
 \le \sqrt\pi L\,\|\dd u^*\|_2
       +2\pi\sum_{T:\,s_T\ne b}r_TA_u(T).
\end{equation}
Consequently, along diagonal transfers $F_n$ with errors tending to zero,
write $u_n^*$ for the transferred coordinate and $D_{n,b}$ for its
distinguished target disk.  Then
\begin{equation}\label{eq:stream-stokes-asymptotic}
 \limsup_{n\to\infty}
 \left|\int_{F_n(C)}\dist(w,D_{n,b})\,\dd u_n^*\right|
 \le \sqrt\pi L\,\|\dd u\|_{2,G\cap U}
       +2\pi L\,\|a_u|_{B\setminus\{b\}}\|_{\ell^2}.
\end{equation}
\end{lemma}

\begin{proof}
Put $\phi(w)=\dist(w,D_b)$.  Its weak gradient has norm at most one.
Apply Stokes' theorem in the region inside $F(C)$ after replacing each
circular hole by a slightly larger concentric circle.  There are only
finitely many holes, so the offsets can be chosen disjoint and inside
the domain.  The exterior boundary integral is the wedge integral
$\int\dd\phi\wedge\dd u^*$ plus the integrals around the offset
circles, each with positive disk orientation.

On the distinguished offset of radius $r_b+\rho$, $\phi$ is the
constant $\rho$.  The period of $u^*$ there is one, so its contribution
is $\rho$ and tends to zero.

At every other hole the period is zero, so there is a single-valued
branch.  Subtracting a constant from that branch and integrating by
parts bounds the boundary integral by
\[
 2\pi(r_T+\rho)\bigl(A_u(T)+o(1)\bigr).
\]
Here $|\dd\phi|\le|\dd w|$ along the Euclidean offset circle, and
the oscillation on small offsets tends
to at most $A_u(T)$ by the collar correspondence in
Lemma~\ref{lem:end-components}.  This argument also covers a hole of
radius zero.  The remaining region has Euclidean area at most
$\pi L^2$, so Cauchy--Schwarz bounds the wedge integral by
$\sqrt\pi L\|\dd u^*\|_2$.  Letting the offsets
tend to zero proves Equation~\eqref{eq:stream-stokes}.

Disjointness of the disk interiors gives $\sum_T r_T^2\le L^2$.
The terminal indices are distinct, and the index $b$ occurs only at
the distinguished hole.  The positive-excess estimate of
Theorem~\ref{thm:transfer} therefore gives
\[
 \left(\sum_{T:\,s_T\ne b}A_u(T)^2\right)^{1/2}
 \le \|a_u|_{B\setminus\{b\}}\|_{\ell^2}+o(1).
\]
Conformal invariance of Dirichlet energy and the energy estimate of
that theorem also give
$\|\dd u^*\|_2\le\|\dd u\|_{2,G\cap U}+o(1)$.
A second application of Cauchy--Schwarz proves
Equation~\eqref{eq:stream-stokes-asymptotic}.
\end{proof}

\subsection{Excluding a protrusion}

We now apply the barrier alternative and the Stokes estimate to a fixed
normalized map and a round disk properly contained in one of its marked
components.  The estimates will exclude a whole neighborhood of that
pair, allowing the final selection to use only countably many tests.

For a fixed mark $b$, let $\mathcal Q_b$ be the space of pairs $(f,D)$
where $f$ is a normalized univalent map on $G$ and $D$ is a closed
round disk, possibly a singleton, contained in $K_b(f)$.  Give maps
the compact-open topology and disks their center-radius topology,
equivalently the Hausdorff topology on the bounded disks considered
here.  These pair spaces are second countable.  A pair is \emph{bad}
if $D\ne K_b(f)$.

\begin{lemma}[Local exclusion of a bad pair]\label{lem:bad-pair}
Given a countable prior list satisfying (P), a mark $b$, and a bad
pair $(f_0,D_0)\in\mathcal Q_b$, one can append a compatible real
barrier or a compatible period coordinate singular at $b$ such that
an entire relative neighborhood of $(f_0,D_0)$ is excluded as limiting
pair data of diagonal finite transfers.  The exclusion persists after
any further countable extensions satisfying (P), provided $b$ remains
marked.
\end{lemma}

\begin{proof}
We first fix geometry that will serve both alternatives.  In the
barrier case it gives a family of line segments with a fixed positive
transverse width.  In the law case it keeps the arc starts a fixed
positive distance from $D_0$.

\smallskip\noindent\emph{Fixed geometry near the protrusion.}
Write $K=K_b(f_0)$.  There is a boundary point $x$ of $K$ outside
$D_0$: for example, maximize distance to the center of $D_0$ on $K$.
Choose a bounded open ball $W_2$ about $x$ whose closure misses $D_0$,
and a smaller closed ball $W_1$ about $x$ contained in $W_2$.
Boundary bumping, applied in a still smaller closed ball, gives a
nondegenerate subcontinuum $Q\subset K\cap\operatorname{int}W_1$.
The ball can be chosen small enough that $K$ also has points outside
it, because $D_0\subset K$ is separated from $x$.

Every boundary point of a complementary component is approached from
the domain.  Indeed, if a ball about such a point missed the domain,
the connected ball would lie in that same complementary component,
contradicting that the point was a boundary point.  Thus choose a
small closed ball $A\subset f_0(G)\cap\operatorname{int}W_1$ and
nested smaller balls strictly inside $A$.  Their preimages give an
anchor $E$ and a compact testing set
$E_*\subset\operatorname{int}E$.  Start the collar sequence deep
enough that its closed disks miss $E$.

There is a direction and a closed interval $J$ of positive transverse
width such that every line with transverse parameter in $J$ joins a
point strictly inside $f_0(E_*)$ to a point of $Q$, by a segment in
$\operatorname{int}W_1$.  To see this, first direct a line from the
center of the smallest anchor ball to a point of $Q$.  That point
projects strictly inside the ball's transverse projection.  Perturb
the direction, if necessary, so that the transverse projection of
$Q$ is nondegenerate.  The projection is an interval, and it overlaps
the interior of the anchor-ball projection in positive length.  Take
$J$ compactly inside that overlap.  Choose the starts on the middle
chords of the anchor ball; they then have a uniform interior buffer.
Convexity of $W_1$ puts the segments inside it.

These sets and $J$ are fixed before the test's collar and penalty
region are chosen.  For every prescribed collar $C_l$, the compact
continuum
$Q$ lies strictly inside the bounded side of $f_0(C_l)$, whereas the
buffered starts lie outside.  After a collar is fixed these inclusions
persist under sufficiently small changes of the map, by winding
number on $C_l$ and on the boundary of $E_*$.  In particular the
buffered starting points remain images of points in $E_*$.

\smallskip\noindent\emph{The barrier case.}
Use the endpoint problem with these anchors and with gates prescribed
by $(f_0,W_1,W_2)$.  Let $c_L$ be a constant large enough for the
line-payment estimate in the common target size bound.  Fix
$\tau>0$ with $c_L\tau<|J|/2$.
If there is a compatible barrier $u$ with
\[
 \|\dd u\|_{2,G\setminus O}+\|a_u\|_{\ell^2}<\tau,
\]
append it.  It equals zero on $E_*$ and one on some fixed inner
collar $C_l$.  In sufficiently close finite-map data the segments just
constructed start in the zero region and reach the side inside
$F(C_l)$; stop each segment at its first hit of that curve.  The
compact set $\overline O$ has image disjoint from $W_1$ under
$f_0$, and this separation persists under convergence on
$\overline O$.  Thus all gradients paid along these segments lie off
the image of $O$.

Lemma~\ref{lem:line-payment}, with the transfer error tending to zero,
gives
\[
 |J|\le c_L\bigl(\|\dd u\|_{2,G\setminus O}
                          +\|a_u\|_{\ell^2}\bigr)+o(1),
\]
a contradiction to $c_L\tau<|J|/2$.  To justify the restricted
energy estimate, retain $\overline O$ in the unchanged
core and use the original energy off $O$ there, together with the
transfer error on the rest.  All required separations involve fixed
compact sets after $u$, $C_l$, and $O$ have been chosen.  This excludes
a relative neighborhood of $f_0$, and hence a relative neighborhood
of $(f_0,D_0)$.

\smallskip\noindent\emph{The law case.}
If no such barrier is available, Theorem~\ref{thm:positive-law}
provides an endpoint law $\mu$ obeying all prescribed gates.  We will
make the Stokes upper bound small while keeping the boundary integral
bounded below.  Put
\[
 d_0=\dist(\overline W_2,D_0)>0.
\]
Apply Lemma~\ref{lem:compact-arcs} at a sufficiently deep collar
$C=\partial U$, and then Lemma~\ref{lem:stream}.  The resulting smooth
period coordinate $u$ can be chosen so that
\begin{equation}\label{eq:bad-pair-small-stream}
 \sqrt\pi L\,\|\dd u\|_{2,G\cap U}
  +2\pi L\,\|a_u|_{B\setminus\{b\}}\|_{\ell^2}
 <d_0/4.
\end{equation}
Indeed the energy norm and the nonsingular budget norm tend to zero
with the depth.  The singular budget is one and is absent from the
Stokes upper bound.
The gate puts all arc starts in $f_0^{-1}(\overline W_2)$, so the
unsmoothed coordinate satisfies
\begin{equation}\label{eq:bad-pair-lower-bound}
 \int_C\dist(f_0(\cdot),D_0)\,\dd v
 =\int\dist(f_0(\alpha(0)),D_0)\,\dd\eta(\alpha)
 \ge d_0.
\end{equation}

We next preserve this lower bound while smoothing.  On the smooth
compact curve $C$, the function $\psi=\dist(f_0(\cdot),D_0)$ is Lipschitz.  Since
$u-v$ is single-valued on $C$, integration by parts for a continuous
function of bounded variation gives
\[
 \left|\int_C\psi\,\dd(u-v)\right|
   =\left|\int_C(u-v)\,\dd\psi\right|
   \le \sup_C|u-v|\,\operatorname{Var}_C(\psi).
\]
After the collar is fixed, choose the smoothing accuracy so fine that
this error is less than $d_0/4$, while maintaining
Equation~\eqref{eq:bad-pair-small-stream}.  Thus
\[
 \int_C\dist(f_0(\cdot),D_0)\,\dd u>3d_0/4.
\]
For fixed smooth $u|_C$, this integral varies continuously with the
map uniformly on $C$ and with the disk in Hausdorff distance.
Consequently it remains greater than $d_0/2$ on a relative neighborhood
$\mathcal V$ of $(f_0,D_0)$.

Append $u$ and retain $b$ as a mark.  In a finite transfer, the
transferred coordinate agrees with $u$ on $C$ for all sufficiently
late stages.  If the limiting pair belonged to $\mathcal V$, the
outer integral in Lemma~\ref{lem:stream-stokes} would therefore have
liminf greater than $d_0/2$.  Equations
\eqref{eq:stream-stokes-asymptotic} and
\eqref{eq:bad-pair-small-stream} give a limsup less than $d_0/4$.
This contradiction excludes $\mathcal V$.

\smallskip\noindent\emph{Persistence.}
In either case the contradiction uses one fixed appended coordinate,
fixed compact source sets, its fixed energy and band bounds, and the
vanishing errors for its eventual transfers.  Enlarging the coordinate
list or the flag set does not change these data.  Later transfers
still satisfy all the estimates once this coordinate and the required
mark are included.  The exclusion therefore persists under all
further compatible countable extensions.
\end{proof}

\subsection{Selection with prescribed retained tests}

The two selection arguments can start with coordinates already supplied
by an application.  We record this stronger conclusion explicitly: the
finite-transfer estimates for those coordinates survive every later
selection, and the limit still has a circle-domain image.

\begin{theorem}[Retained-test transfer and selection]\label{thm:retained-tests}
Use the quotient sphere $X=G\cup B$, fixed exterior radius $R$, and
coordinates of Definition~\ref{def:coordinates}.  Let $\mathscr L_0$
be any countable well-ordered list satisfying compatibility~\textup{(P)},
with a countable flag set $\Sigma_0$.  Prescribe also a countable set
$S_0\subset B$ of mandatory marks containing every singularity of
$\mathscr L_0$.

There are a countable well-ordered compatible extension $\mathscr L$
of $\mathscr L_0$, with enlarged countable flags, and a countable mark
set $S\supset S_0$ containing all its singularities, with the following
properties.  Every normalized diagonal finite-transfer limit of
$\mathscr L$, retaining the marks $S$, has a circle-domain image.
This conclusion persists after any further countable compatible
extension, retaining $S$ and including all additional singularities among
the marks.

More precisely, let $\varepsilon_n>0$ tend to zero and let
$H_n\subset G$ be increasing compact sets such that every compact subset
of $G$ is contained in $\operatorname{int}H_n$ for all sufficiently
large $n$.  One can choose increasing finite subsets $I_n$ exhausting
the labels of $\mathscr L$, finite bordered completions $(Q_n,g_n)$,
nested smooth cores $K_n\subset G$, transferred coordinates $u_{i,n}$,
and normalized conformal maps
\[
 f_n:Q_n^\circ\longrightarrow\Omega_n,
 \qquad f_n(z)=z+O(1/z),
\]
where $\Omega_n$ is a finite circle domain, such that $H_n\subset
\operatorname{int}K_n$ and the following assertions hold.
\begin{enumerate}[label=\textup{(\roman*)}]
\item The metric $g_n$ is the original spherical metric near $K_n$.
For every $i\in I_n$, the coordinate $u_{i,n}$ agrees with $u_i$ near
$K_n\cap M_i$, and all its terminal periods are those specified in
Theorem~\ref{thm:transfer}.
\item The total additional energy satisfies
\begin{equation}\label{eq:retained-total-tail}
 \sum_{i\in I_n}
 \|\dd u_{i,n}\|_{2,(Q_n\setminus K_n)\cap M_i,g_n}^2
 <\varepsilon_n^2.
\end{equation}
For every $i\in I_n$, if $T_{n,s}\subset M_i$ is a terminal disk
and $s$ is nonsingular for $u_i$,
its boundary oscillation $A_{i,n}(T_{n,s})$ satisfies
\begin{equation}\label{eq:retained-excess-oscillation}
 \left(\sum_{\substack{s:\ T_{n,s}\subset M_i\\s\ne p_i}}
 [A_{i,n}(T_{n,s})-a_i(s)]_+^2\right)^{1/2}
 <\varepsilon_n.
\end{equation}
The singularity exclusion is omitted for real coordinates.
\item Every prescribed mark eventually indexes a terminal disk.
All terminal disks have quotient diameter tending uniformly to zero.
After passing to a subsequence, $f_n\to f$ locally uniformly on $G$,
where $f$ is normalized and univalent and $f(G)$ is a circle domain.
At each marked end, its distinguished target disks converge to the
entire component $K_b(f)$.
\item All complementary disks of $\Omega_n$, and all images of
physical points in $Q_n^\circ$ with $|z|\le4R$, lie in $|w|<5R$.  Any prescribed
compact source sets, including finitely many smooth contours and
their collars at each stage, can be retained in $K_n$.  A retained
oriented dust contour has the bounded-side correspondence of
Lemma~\ref{lem:end-components}.  For each fixed such contour its
bounded target sides lie in a fixed disk, uniformly in $n$ and in
the number of retained tests.
\end{enumerate}
The finite subsets $I_n$ carry their inherited order from
$\mathscr L$; they need not be initial segments of its well-order.
The original coordinates, their masks, and their budgets are unchanged.
\end{theorem}

\begin{proof}
\emph{First selection with the seed retained.}
Start the selection in Section~\ref{sec:selection} with
$\mathscr L_0$ and $\Sigma_0$.  Every application of
Theorem~\ref{thm:countable-exceptions} permits precisely this arbitrary
current countable list and permits a countable flag enlargement.
At a successor stage append the selected barrier after the existing
well-order.  At a countable limit stage take the order-preserving
union of all earlier lists and the union of the flags.  These unions
are countable.  Each compatibility requirement concerns one new
label and finitely many lower labels, so it already holds at a
previous stage; increasing the flags can only remove free points.
Thus compatibility persists, while all previously chosen functions
and their bounds remain fixed.

The countable-basis argument in Section~\ref{sec:selection} makes
this process stop at a countable ordinal.  It produces a countable
set $S_1$ as in Equation~\eqref{eq:first-exceptional-set}, with
$K_b(f)$ a singleton for every $b\notin S_1$ in every later
compatible transfer limit.  Put $S=S_0\cup S_1$.  This enlargement
is countable and contains the seed singularities.  The newly selected
barriers are real and introduce no singularities.  In particular,
Theorem~\ref{thm:unmarked-collapse} applies outside $S$.

\emph{Second selection at every mandatory mark.}
Consider the second-countable disjoint union
\[
 \mathcal Q=\coprod_{b\in S}\mathcal Q_b.
\]
Starting with an empty excluded open subset of $\mathcal Q$, select
an uncovered bad pair whenever one exists.  Apply
Lemma~\ref{lem:bad-pair}, append its coordinate, and add its excluded
relative neighborhood to the open set.  At countable limit stages
take the order-preserving unions of coordinate lists and flags.
All newly introduced period singularities belong to $S$.

This selection also stops at a countable stage.  Each strict
increase of the excluded open set contains a basic open set not
contained in the preceding one.  Once chosen, that basic set is
contained in all later excluded open sets, so it cannot be chosen
again.  An $\omega_1$-long selection would contradict the countability
of the basis.  At termination every bad pair is covered, and the
accumulated list and flags are countable.  The exclusions persist by
Lemma~\ref{lem:bad-pair}.

\emph{Finite transfers and preservation of the original tests.}
Enumerate the underlying countable set of final labels and take
increasing finite subsets $I_n$ exhausting it, each ordered by the
inherited well-order.  A countable well-order need not have finite
initial segments exhausting it; none are needed here.
Similarly choose increasing finite subsets exhausting $S$ and include
in the $n$th prescribed terminal indices every singularity of $I_n$.
Apply Theorem~\ref{thm:transfer} to $I_n$, with error
$\varepsilon_n$, prescribed dust-cover mesh tending to zero, and
input core containing $H_n$, the preceding core, and all additionally
specified compact source sets.  One may include compact neighborhoods
of these sets, so the resulting cores are nested and contain $H_n$
in their interiors.  Include the fixed exterior region and all current
mask boundaries as required by that theorem.

Finite circle-domain uniformization gives $f_n$ and $\Omega_n$.
Exact core agreement, the total tail bound, and the individual
positive-excess bounds are exactly the conclusions of
Theorem~\ref{thm:transfer}.  They apply to every fixed seed coordinate
for all sufficiently large $n$.  No selected flag enlargement changes
a coordinate's original band budget.  The terminal disks shrink in
the quotient metric by the prescribed cover meshes.  Normalized
compactness and the uniform radius bound follow from
Lemma~\ref{lem:normalized-limits}.

\emph{Identification of the limiting components.}
Pass to a subsequence on which the maps and all marked disks converge.
Lemma~\ref{lem:end-components} gives $D_b\subset K_b(f)$ for
$b\in S$.  If one inclusion were strict, $(f,D_b)$ would be a bad
pair in a previously excluded neighborhood.  The fixed coordinate
responsible for that neighborhood occurs in every sufficiently late
finite sublist.  Its exact core values and vanishing transfer errors
contradict its exclusion estimate.  Thus $D_b=K_b(f)$ for every
$b\in S$.  All components outside $S$ are points by
Theorem~\ref{thm:unmarked-collapse}, and these are all components by
Lemma~\ref{lem:end-components}.  Hence $f(G)$ is a circle domain.

Every argument just used depends on fixed retained coordinates and
marks and survives further compatible extensions.  A new marked end
outside $S$ already has a singleton component; its distinguished disk
is contained in that singleton and therefore equals it.  This also
proves the stated persistence for additional singularities and marks.

Finally, retain a fixed contour and its collar in all sufficiently
late cores.  Orientation and Jordan separation give its bounded-side
correspondence in each finite surface.  For the bound on its image,
choose $r_C>2R$ so large that the source contour is inside $|z|<r_C$.
The exterior estimate in Equation~\eqref{eq:transfer-exterior-bound}
bounds $f_n(\{|z|=r_C\})$ uniformly; the contour's image and its
bounded side lie inside that outer image curve.  This gives a bound
depending only on $R$ and the fixed contour, independent of the number
of tests.  Contours inside $|z|\le4R$ have the common bound $5R$.
\end{proof}

\begin{proof}[Proof of Theorem~\ref{thm:main}]
For a proper sphere domain make an initial M\"obius change of source
coordinate so that infinity belongs to the domain.  Apply
Theorem~\ref{thm:retained-tests} with the empty seed list and no
prescribed marks.  It supplies a conformal circle-domain image.
The whole-sphere case is immediate.
\end{proof}

The following form uses continuous real tests on the quotient sphere.
It is the retained-test input used in the rigidity companion.  Its
trace hypothesis is especially simple: every pair of boundary values
has a planar image of area zero.

\begin{corollary}[Retained scalar tests]\label{cor:scalar-retained-tests}
Let $G\subset\Chat$ contain infinity, let $X=G\cup B$ be its quotient
sphere, and fix $R$ with $\{|z|\ge R\}\cup\{\infty\}\subset G$.
Let $u_1,\ldots,u_m\in C(X,\R)$ be smooth on $G$, with
$\int_G|\dd u_i|^2<\infty$, and suppose
\begin{equation}\label{eq:scalar-pair-null}
 \area\bigl((u_i,u_j)(B)\bigr)=0\qquad(i\ne j).
\end{equation}
For every positive sequence $\varepsilon_n\to0$ there are finite
bordered completions $(Q_n,g_n)$, nested smooth cores $K_n\subset G$
containing every compact subset of $G$ in their interiors eventually,
normalized conformal maps $f_n:Q_n^\circ\to\Omega_n$ onto finite
circle domains, and smooth real tests $v_{i,n}$ on $\Omega_n$ such that
\begin{align}
 v_{i,n}\circ f_n&=u_i
   &&\text{near }K_n,\label{eq:scalar-core}\\
 \sum_{i=1}^m\int_{\Omega_n\setminus f_n(K_n)}
             |\nabla v_{i,n}|^2\,\dd a
   &<\varepsilon_n^2,\label{eq:scalar-tail}\\
 \sum_T A_{i,n}(T)^2&<\varepsilon_n^2
   &&(1\le i\le m).\label{eq:scalar-oscillation}
\end{align}
Here $T$ ranges over the complementary round disks of $\Omega_n$;
the cluster values of $v_{i,n}$ at $T$ lie in an interval of length
$A_{i,n}(T)$.  The metric $g_n$ is the original metric near $K_n$.
After passage to a subsequence, $f_n\to f$ locally uniformly on $G$,
where $f$ is normalized and univalent and $f(G)$ is a circle domain.
All finite target holes lie in $|w|<5R$.  The source-contour retention,
correspondence, and bounds of Theorem~\ref{thm:retained-tests}\textup{(iv)}
also hold.  In particular, these conclusions allow arbitrarily many,
and possibly uncountably many, point components in the limiting
complement.
\end{corollary}

\begin{proof}
If $B=\varnothing$, take $Q_n=K_n=G=\Chat$, $f_n$ the identity,
and $v_{i,n}=u_i$; all tail and oscillation sums vanish.  We may
therefore assume that $B$ is nonempty.

Each test is bounded by compactness of $X$, has full mask, and has
singleton band $\{u_i(b)\}$ at every $b\in B$.  It is therefore a
coordinate with budget $a_i=0$, and the initial flag set may be empty.
We verify compatibility in any chosen ordering of the finite list.

For distinct labels $i,j$, Equation~\eqref{eq:scalar-pair-null}
implies that the image of $B$ under $(u_i,u_j)$ modulo one has zero
Haar measure on $\T^2$: decompose $\R^2$ into countably many integer
translates of a unit square.  For a label $j$ and a finite tuple of
lower labels $i_1,\ldots,i_d$, a common dust point between $H_j(t)$
and any $H_{i_k}(t_k)$ requires $(t,t_k)$ to lie in this null pair
image.  Each such condition is a null cylinder in $\T^{d+1}$, by
Fubini.  Their finite union is null, including when lower labels are
repeated.  Almost every tuple therefore has no such common dust
point at all, which is stronger than compatibility~\textup{(P)}.

Apply Theorem~\ref{thm:retained-tests} with these seeds, choosing every
finite sublist to contain all $m$ seeds from the outset.  Put
$v_{i,n}=u_{i,n}\circ f_n^{-1}$.  Exact core agreement gives
Equation~\eqref{eq:scalar-core}, and conformal invariance of energy
gives Equation~\eqref{eq:scalar-tail}.  Since all seed budgets vanish,
Equation~\eqref{eq:retained-excess-oscillation} gives
Equation~\eqref{eq:scalar-oscillation}.  The cluster-interval assertion
is Lemma~\ref{lem:end-components}.  All remaining conclusions are
those of the theorem.
\end{proof}
